\chapter*{IDENTITAS DAN URAIAN UMUM PENELITIAN}
\addcontentsline{toc}{chapter}{IDENTITAS DAN URAIAN UMUM PENELITIAN}

\begin{enumerate}[itemsep=0.2cm]
\item Judul Penelitian: \judulPenelitian

\item Tim Peneliti:

\vspace{0.1cm}
{\small
\begin{longtable}{|c|p{3cm}|p{1.8cm}|p{2.5cm}|p{2.5cm}|p{1.3cm}|}
\hline
\textbf{No} & \textbf{Nama} & \textbf{Jabatan} & \textbf{Bidang Keahlian} & \textbf{Jurusan} & \textbf{Alokasi Waktu} \\
\hline
\endfirsthead
\hline
\textbf{No} & \textbf{Nama} & \textbf{Jabatan} & \textbf{Bidang Keahlian} & \textbf{Jurusan} & \textbf{Alokasi Waktu} \\
\hline
\endhead
1 & \ketuaNama\newline NIDN: \ketuaNIDN & Ketua & Rekayasa Perangkat Lunak & Teknik Informatika & 12 \\
\hline
2 & \anggotaSatuNama\newline NIDN: \anggotaSatuNIDN & Anggota 1 & Basis Data dan Sistem Informasi & Teknik Informatika & 10 \\
\hline
3 & \anggotaDuaNama\newline NIDN: \anggotaDuaNIDN & Anggota 2 & Kecerdasan Buatan dan Data Science & Teknik Informatika & 10 \\
\hline
4 & \anggotaTigaNama\newline NIDN: \anggotaTigaNIDN & Anggota 3 & Jaringan Komputer dan Sistem Terdistribusi & Teknik Informatika & 10 \\
\hline
5 & \mahasiswaSatuNama\newline NIM: \mahasiswaSatuNIM & Pembantu pelaksana & Teknik Informatika & Teknik Informatika & 8 \\
\hline
6 & \mahasiswaDuaNama\newline NIM: \mahasiswaDuaNIM & Pembantu pelaksana & Teknik Informatika & Teknik Informatika & 8 \\
\hline
7 & \mahasiswaTigaNama\newline NIM: \mahasiswaTigaNIM & Pembantu pelaksana & Teknik Informatika & Teknik Informatika & 8 \\
\hline
8 & \mahasiswaEmpatNama\newline NIM: \mahasiswaEmpatNIM & Pembantu pelaksana & Teknik Informatika & Teknik Informatika & 8 \\
\hline
\end{longtable}
}

\item Objek Penelitian: Aplikasi berbasis web menggunakan dataset internal Jurusan Teknologi Informasi sebagai studi kasus.

\item Masa Pelaksanaan: \lamaPenelitian\ bulan (\bulanTahun)

\item Biaya yang Diusulkan: Rp \biayaDiusulkan,-

\item Lokasi Penelitian: Jurusan Teknologi Informasi, Politeknik Negeri Malang

\item Instansi Lain yang Terlibat: Tidak ada

\item Temuan yang Ditargetkan (TKT): Prototipe TathyaTest dengan target TKT 4, yaitu validasi komponen dalam lingkungan laboratorium. TathyaTest sudah dapat digunakan untuk pengujian aplikasi web dengan dataset internal, meskipun masih dalam skala terbatas.

\item Kontribusi Mendasar pada Suatu Bidang Ilmu: Penelitian ini memberikan kontribusi pada bidang Software Testing dan Quality Assurance melalui pengembangan metode otomatis untuk test case generation yang memanfaatkan web crawling dan ekstraksi elemen HTML. TathyaTest mengotomatisasi proses ini dengan memanfaatkan dataset internal untuk menghasilkan test case yang lebih relevan dan kontekstual.

\item Jurnal Ilmiah yang Menjadi Sasaran: Jurnal nasional terakreditasi Sinta 3 Sisfokom pada bidang Teknologi Informasi.

\item Rencana Luaran HKI: Hak Cipta untuk prototipe TathyaTest dengan status Granted.

\end{enumerate}
